\section{Introduction}


\subsection{Background}

Shor's algorithm~\cite{shor1999} renders RSA, ECDSA, and all
discrete-logarithm-based signature schemes insecure on a sufficiently
powerful quantum computer. NIST's Post-Quantum Cryptography
standardization process~\cite{nist-pqc} selected ML-KEM (Kyber),
ML-DSA (Dilithium), and SLH-DSA (SPHINCS+) as the primary
replacements. Lux Network adopted these standards at the consensus
layer and extended them with ring-signature privacy (Corona) and
fully homomorphic encryption (TFHE) for computation over ciphertexts.

\subsection{Scope}

This paper covers five layers of Lux's post-quantum stack:

\begin{enumerate}[nosep]
  \item \textbf{Consensus signatures}: BLS-381 + ML-DSA-65 + Corona
        (Quasar protocol, LP-020).
  \item \textbf{Threshold key management}: TFHE key generation via
        Shamir LSSS over a 2048-bit safe prime, with Feldman VSS
        commitments.
  \item \textbf{Fully homomorphic encryption}: TFHE boolean-gate
        circuits for encrypted CRDT merge, encrypted search, and
        private smart contracts.
  \item \textbf{Non-custodial wallets}: 2-of-3 MPC (CGGMP21 for
        secp256k1, FROST for Ed25519) with HSM co-signing.
  \item \textbf{Regulatory disclosure}: viewing keys, threshold
        disclosure, and ZK attestation anchoring.
\end{enumerate}

\subsection{Evolution from the 2022 Design}

The 2022 internal document proposed Lamport one-time signatures for
vault protection and an OTP-based information-theoretic signing layer.
Both were superseded:

\begin{itemize}[nosep]
  \item Lamport signatures $\rightarrow$ ML-DSA-65 (FIPS~204): multi-use,
        smaller signatures (3.3\,KB vs 16\,KB per Lamport instance),
        NIST-standardized.
  \item Network-layer OTP signing $\rightarrow$ dropped. The throughput
        cost of per-transaction OTP generation was incompatible with
        sub-second finality. ML-DSA provides computational (not
        information-theoretic) PQ security, which is sufficient given
        current lattice-hardness estimates~\cite{nist-pqc-report}.
\end{itemize}

%% ============================================================
